\chapter{Transmission Lines, Reflection, and SWR}
\section{Why Lines Matter}
Transmission lines are distributed systems. When a feed line is a significant
fraction of a wavelength, voltage and current vary along the line and a lumped
wire approximation fails.

\begin{controlsbox}
A transmission line is closer in spirit to a distributed-parameter dynamical
system than to a lumped ODE. Wavelength matters because position becomes part of
the state description.
\end{controlsbox}

\section{Wavelength and Velocity Factor}
\[
\lambda=\frac{v_p}{f}=\frac{VF\,c}{f}.
\]
Velocity factor \(VF\) is the ratio of wave speed in the line to the speed of
light in vacuum. Physical length is shorter than free-space electrical length
when \(VF<1\).

\section{Characteristic Impedance}
Characteristic impedance \(Z_0\) is the voltage-to-current ratio of a traveling
wave on the line. Common coaxial systems use \(\SI{50}{\ohm}\), but other values
such as \(\SI{75}{\ohm}\), \(\SI{300}{\ohm}\), and \(\SI{450}{\ohm}\) are also
common in specific applications.

\section{Reflection Coefficient}
At the load,
\[
\Gamma=\frac{Z_L-Z_0}{Z_L+Z_0}.
\]
If \(Z_L=Z_0\), then \(\Gamma=0\) and there is no reflection. If the load is open
or shorted, \(|\Gamma|=1\) in the lossless idealization.

\section{SWR and Return Loss}
\[
\mathrm{SWR}=\frac{1+|\Gamma|}{1-|\Gamma|}.
\]
SWR equals 1 for a perfect match and increases as mismatch worsens. Return loss
is
\[
RL=-20\log_{10}|\Gamma|.
\]
Higher return loss means a smaller reflection.

\section{Quarter-Wave Transformation}
For a lossless quarter-wave section,
\[
Z_{\mathrm{in}}=\frac{Z_t^2}{Z_L}.
\]
To match a real load \(Z_L\) to a real line impedance \(Z_0\),
\[
Z_t=\sqrt{Z_0Z_L}.
\]
A half-wave line repeats the load impedance at its input, in the lossless ideal
case.

\section{Smith Chart Interpretation}
A Smith chart is a graphical map of normalized complex impedance and reflection
coefficient. It helps visualize resistance circles, reactance arcs, constant
SWR circles, and how impedance changes as one moves along a transmission line.
For this book, the key point is conceptual: the chart is not a new law of
circuits; it is a coordinate transform for impedance and reflection.

\begin{figure}[htbp]
\centering
\begin{tikzpicture}[scale=2.0]
\draw[thick] (0,0) circle (1);
\draw[GuideBlue,thick] (-1,0)--(1,0);
\draw[GuideLine] (0.333,0) circle (0.667);
\draw[GuideLine] (0.6,0) circle (0.4);
\draw[GuideGreen] (1,0) arc (0:180:0.5);
\draw[GuideGreen] (1,0) arc (0:-180:0.5);
\fill (0,0) circle (0.015) node[below left]{match};
\node[below] at (0,-1.12){normalized impedance plane mapped to \(\Gamma\)};
\end{tikzpicture}
\caption{Simplified Smith-chart geometry: constant-resistance circles and reactance arcs.}
\end{figure}

\section{Worked Example}
A \(\SI{100}{\ohm}\) load terminates a \(\SI{50}{\ohm}\) line:
\[
\Gamma=\frac{100-50}{100+50}=\frac{1}{3}.
\]
\[
\mathrm{SWR}=\frac{1+1/3}{1-1/3}=2.
\]
The return loss is
\[
RL=-20\log_{10}(1/3)\approx \SI{9.54}{\decibel}.
\]
